\chapter{Coated Tongue}

\section*{Definition}
A coated tongue (furred tongue) is the accumulation of desquamated epithelial cells, bacteria, food debris, and salivary proteins on the dorsal tongue surface, appearing as a white or yellowish-white deposit that may be localized or diffuse.

\section*{When to See a Doctor}
See your doctor if:
\begin{itemize}
  \item Coated tongue persists > 2--3 weeks despite improved oral hygiene
  \item A persistent white, red, black, or hairy-appearing patch, especially if it cannot be gently removed
  \item Pain, burning, or bleeding of the tongue
  \item Associated dysgeusia (altered taste) or halitosis
  \item Systemic symptoms (fever, weight loss, dysphagia) accompanying the coat
\end{itemize}

\section*{How It Develops}
Normal tongue desquamation is balanced by mechanical debridement from chewing, saliva flow, and oral microflora. When oral hygiene decreases, salivary flow diminishes, or diet changes to soft foods, the filiform papillae elongate and trap debris, forming the visible coat.

\section*{Causes}
\begin{itemize}
  \item \textbf{Poor oral hygiene:} Inadequate tongue brushing or overall dental care
  \item \textbf{Xerostomia:} Dehydration, medications (anticholinergics, diuretics), Sjogren syndrome, radiation therapy
  \item \textbf{Soft food diet:} Liquid or pureed diets lacking abrasive mechanical cleaning
  \item \textbf{Febrile illness:} Dehydration and reduced oral intake during acute infection
  \item \textbf{Tobacco use:} Smoking or chewing tobacco alters oral flora and reduces papillae desquamation
  \item \textbf{Oral candidiasis:} White plaques that may be confused with benign coated tongue; scrapable
  \item \textbf{Antibiotic use:} Disruption of normal oral microbiome
\end{itemize}

\section*{Related Symptoms}
\begin{itemize}
  \item Halitosis (bad breath)
  \item Dysgeusia (altered taste, metallic or bitter)
  \item Xerostomia (dry mouth)
  \item White or yellowish dorsal tongue discoloration
  \item Loss of taste acuity
\end{itemize}
\textbf{Important:} A simple coating is often benign and reversible, but appearance alone cannot reliably distinguish it from candidiasis, leukoplakia, oral hairy leukoplakia, or another lesion.

\section*{Medical Evaluation}
\begin{itemize}
  \item Inspection of the tongue dorsum: color, thickness, distribution, and scrapability of the coat
  \item Oral examination for signs of candidiasis (erythema, scrapable plaques), leukoplakia, or lichen planus
  \item Review of medications (anticholinergics, antipsychotics, antibiotics) and medical history (diabetes, immunosuppression)
  \item Salivary-flow testing or fungal testing is reserved for selected cases when dry mouth or candidiasis is suspected
  \item Biopsy or specialist referral may be needed for a persistent unexplained lesion, particularly one that is non-scrapable, indurated, ulcerated, or bleeding
\end{itemize}

\section*{Treatment Options}
\begin{itemize}
  \item Improved tongue hygiene (scraping or brushing) -- usually resolves within days to weeks
  \item Treatment of dry mouth with hydration, saliva substitutes, and clinician-directed therapy when appropriate
  \item Nystatin suspension (swish and swallow) or clotrimazole troches for confirmed candidiasis
  \item Discontinuation of causative medication if feasible (anticholinergics)
  \item Smoking cessation counseling
  \item Dietary modification to include abrasive, fibrous foods
  \item Reassurance: benign coated tongue requires no further treatment
\end{itemize}

\section*{Self-Care and Home Management}
\begin{itemize}
  \item Gentle tongue scraping with a dedicated tongue scraper or the edge of a spoon once daily
  \item Brushing the tongue dorsum with a soft-bristled toothbrush
  \item Sip water regularly according to thirst and medical advice
  \item Sugar-free gum or lozenges to stimulate saliva flow
  \item Avoid tobacco and alcohol-based mouthwashes
  \item Increase dietary roughage (raw vegetables, fibrous fruits)
  \item Mild saline mouth rinse (1/2 tsp salt in 8 oz warm water) once daily
\end{itemize}

\section*{References}
\begin{thebibliography}{99}
\bibitem{coated_tongue:ref1} Byrd JA, Bruce AJ, Rogers RS III. ``Glossitis and Other Tongue Disorders.'' Dermatol Clin. 2003;21(1):145--154.
\bibitem{coated_tongue:ref2} Reamy BV, Derby R, Bunt CW. ``Common Tongue Conditions in Primary Care.'' Am Fam Physician. 2010;81(5):627--634.
\bibitem{coated_tongue:ref3} Scully C. ``Oral and Maxillofacial Medicine.'' 3rd ed. Churchill Livingstone, 2013.
\bibitem{coated_tongue:ref4} Neville BW, Damm DD, Allen CM, et al. ``Oral and Maxillofacial Pathology.'' 4th ed. Elsevier, 2016.
\bibitem{coated_tongue:ref5} Gurvits GE, Tan A. ``Black Hairy Tongue Syndrome.'' World J Gastroenterol. 2014;20(31):10845--10850.
\end{thebibliography}
